%%%%%%%%%%%%%%%%%%%%%%%%%%%%%%%%%%%%%%%%%%%%%%%%%%%%%%%%%%%%%%%%%%%%%%%%
\section{Open problems}\label{sec:open}
%%%%%%%%%%%%%%%%%%%%%%%%%%%%%%%%%%%%%%%%%%%%%%%%%%%%%%%%%%%%%%%%%%%%%%%%

\begin{enumerate}
\item \textbf{The sharp continuum exponent.} Prove $\Prob(\Pi\in E_\delta)\le\delta^{\astar-o(1)}$ for the
meandric permuton, for instance by showing that a near-split forces a neck of boundary length of
order $\sqrt\delta$, so that the glued-disk calculation of Section~\ref{sec:disks} could finish.
% UNVERIFIED: whether improving the one-pin conditional estimate alone to delta^(1-o(1)) suffices
% for the sharp exponent is not stated in the sources; omitted here.
\item \textbf{A local law in the APPS model.} Prove $\Prob(N=n)=n^{-1-4/\gamma^2+o(1)}$ along
$n\equiv1\pmod3$; the upper bound is proved.
\item \textbf{Determinant-weighted maps.} Relate planar maps weighted by powers of their number of
spanning trees to the determinant-weighted APPS model, with enough control of rare events to
transfer Theorem~\ref{thm:kpz}.
\item \textbf{Convergence.} Prove Conjecture~\ref{conj:bgs}, or even its weaker permuton form,
Conjecture~\ref{conj:permuton}.
\item \textbf{A separation lemma.} Show that near-splits of finite meanders at moderate tolerances,
much coarser than $(\log n)/n$, are only a factor $n^{o(1)}$ more likely than exact ones, a meander
analog of Kesten's estimates. It would allow coarser tolerances in
Proposition~\ref{prop:condlower}.
\item \textbf{A quantitative coupling} as in Proposition~\ref{prop:condupper}, with a polynomial rate.
\end{enumerate}
